\begin{abstract}
Large Language Models can generate REST API tests from an OpenAPI specification, and a
growing literature reports that they find more errors than search-based tools. Surveying
59 papers (2018--2026), we find no study comparing an LLM against both a manual suite and
a mature tool, none computing fault-detection Recall against a known fault population, and
none reporting a per-endpoint edge-case metric.

We compare three arms blind: an LLM (Claude Sonnet 4.6), an independently authored manual
EP/BVA suite, and EvoMaster 6.0.0. The subjects are three EvoMaster Benchmark APIs covering
\RQIopsTotal{} operations, into which we seed \RQIIn{} faults by mutation, fixing the
denominator of Recall exactly.

The LLM attains \RQIcoveragePct\% endpoint coverage and writes \RQIIItotalLlm{} edge-case
scenarios to the manual suite's \RQIIItotalManual{}. Ground-truth Recall nonetheless orders
the arms EvoMaster \RQIIrecallEvo{} $>$ LLM \RQIIrecallLlm{} $>$ manual \RQIIrecallManual{}.
The ordering follows from oracle kind: a specification-derived status oracle cannot see a
fault that changes a computed value while preserving the HTTP status. Scenario counts
therefore mislead when reported without a ground-truth Recall beside them.

\keywords{REST API testing \and Large language models \and Mutation testing
\and Test oracles \and Empirical software engineering}
\end{abstract}
